\documentclass[11pt,a4paper]{article}

\usepackage[margin=2.5cm]{geometry}
\usepackage{amsmath,amssymb,amsthm}
\usepackage{algorithm}
\usepackage{algpseudocode}
\usepackage{booktabs}
\usepackage{graphicx}
\usepackage{xcolor}
\usepackage{hyperref}

\hypersetup{colorlinks=true,linkcolor=blue!50!black,citecolor=blue!50!black,
            urlcolor=blue!50!black}

\newtheorem{proposition}{Proposition}
\newtheorem{remark}{Remark}

\newcommand{\xb}{\mathbf{x}^{b}}
\newcommand{\xa}{\mathbf{x}^{a}}
\newcommand{\Bmat}{\mathbf{B}}
\newcommand{\Lmat}{\mathbf{L}}
\newcommand{\Dmat}{\mathbf{D}}
\newcommand{\Hmat}{\mathbf{H}}
\newcommand{\Rmat}{\mathbf{R}}

\title{\textbf{Adaptive localization by observation assignment\\
for the modified-Cholesky ensemble Kalman filter}}

\author{Working draft}
\date{\today}

\begin{document}
\maketitle
\begin{center}\fbox{\parbox{0.92\linewidth}{\small\textbf{Draft status.} The selection rule of Section~3 (the cost $J$) is superseded: it selects the least correlated candidate (see \texttt{nota\_criterio.tex}). The method now reads each component's radius from the partial correlations of a probe modified-Cholesky precision (\texttt{qgloc/assignment.py}, \texttt{assign\_partial}); the analysis of Section~4 and the benchmark design are unchanged. Sections~3 and~5 are to be rewritten from the tables \texttt{tables\_paper.tex} once EXP-02 has run.}}\end{center}

\begin{abstract}
\noindent
Localization in ensemble data assimilation is normally parameterized by a
radius of influence, tuned once against a known truth and then held fixed. We
replace the radius by a different object: each model component is updated by
exactly one observation, chosen from the few that its neighbourhood reaches,
and the radius is whatever distance that choice implies. The choice is made by
minimizing a variational cost evaluated at a scalar local analysis, computed
from ensemble anomalies and observations alone, with no reference trajectory
and nothing withheld. Because the cost decomposes over components, the optimal
assignment is obtained exactly, in closed form, at a cost of microseconds per
evaluation. The resulting radius field adapts to the local observation density
without any parameter controlling it, and the observation selected is very
often not the nearest one: on a quasi-geostrophic testbed the second, third and
fourth candidates are chosen about as often as the first, because what decides
the choice is the ensemble correlation and not the distance. The assignment
determines the predecessor sets of a modified-Cholesky precision matrix, which
is assembled once per cycle and used for a single global analysis.
\end{abstract}

\section{The parameter nobody can pin down}

The analysis step of an ensemble Kalman filter combines a forecast ensemble
with observations. With an ensemble far smaller than the state dimension, the
sample covariance is rank deficient and carries spurious correlations between
distant variables; localization removes them. Every localization scheme depends
on a radius of influence, and the radius is normally chosen by sweeping it
against a known truth in a synthetic experiment and then freezing it.

Attempts to estimate the radius instead run into a difficulty that is not a
matter of tuning. A radius per model component is a free parameter with nothing
in the data to pin it down. In experiments preceding this work, on a system of
40 components, an oracle with full access to the truth reduced the analysis
error by 59\% on the cycle it optimized, and the same radius field made the
error 22\% \emph{worse} on the following cycle, performing worse than a random
shuffle of its own values. What such a search finds is the cycle's noise.

Grouping components reduces the count but introduces a grouping that must
itself be justified, and the radius within a group remains a free number.

\section{Assignment instead of radius}

We propose a different parameterization. Let the model state have components
$i = 1, \dots, n$ and let there be $m$ observations at known locations. For each
component, grow a neighbourhood until it contains $\kappa$ observations; call
them the \emph{candidates} of $i$ and write $\mathcal{C}(i)$ for the set. The
component is then updated by exactly one member of $\mathcal{C}(i)$, or by none
at all.

The decision variable is the assignment
%
\begin{equation}
  a : \{1, \dots, n\} \;\longrightarrow\; \mathcal{C}(i) \cup \{\varnothing\},
\end{equation}
%
and the radius of component $i$ is not chosen: it is the distance to $a(i)$.

Three properties follow immediately.

\paragraph{The radius adapts to the observation density with no parameter.}
Where observations are dense the neighbourhood stops growing early; where they
are sparse it grows further. On the testbed below, a regular lattice of spacing
four gives a mean radius of 2.7 and a maximum of 4, while a random network of
the same overall density gives a mean of 1.7 and a maximum of 5 --- the second
number is larger precisely because a random network leaves empty regions. No
setting controls this.

\paragraph{The nearest observation is not automatically the right one.}
What determines how much an observation tells us about a component is the
ensemble correlation between them, not the distance. In a flow with structure,
two points aligned along a current can be far better correlated than two
adjacent points across a front. This is why the neighbourhood keeps growing
past the first observation: the choice is made on the cost, and the cost sees
the correlation.

\paragraph{Abstaining is a legal choice.}
If none of the candidates correlates with the component, updating it from any
of them injects noise, which is what localization exists to prevent. Leaving
the component at its background value is one of the options, and the cost
selects it on its own.

\section{The cost}

For component $i$ assigned to observation $j$, the local analysis is the scalar
Kalman update
%
\begin{equation}
  x^{a}_{i} = x^{b}_{i} + \frac{\sigma_{ij}}{\sigma_{jj} + r_{j}}
              \big( y_{j} - x^{b}_{j} \big),
  \label{eq:local}
\end{equation}
%
with $\sigma_{ij}$ the ensemble covariance between component $i$ and the
observed site of $j$, $\sigma_{jj}$ the ensemble variance there, and $r_{j}$
the observation error variance. The cost is the variational one evaluated at
that analysis,
%
\begin{equation}
  J_{i}(a) \;=\;
  \underbrace{\frac{(x^{a}_{i} - x^{b}_{i})^{2}}{\sigma_{ii}}}_{\text{departure from the background}}
  \;+\;
  \underbrace{\frac{(y_{j} - x^{a}_{j})^{2}}{r_{j}}}_{\text{residual}} ,
  \qquad j = a(i),
  \label{eq:cost}
\end{equation}
%
and $J_{i} = (y_{j} - x^{b}_{j})^{2} / r_{j}$ when $i$ abstains, since the
background is kept and the departure term vanishes. The global score is
%
\begin{equation}
  J(a) \;=\; \frac{1}{n}\sum_{i=1}^{n} J_{i}(a).
  \label{eq:global}
\end{equation}

No reference trajectory enters \eqref{eq:cost}, and no observation is withheld.
Everything is a dot product over the $N$ ensemble members. An assignment that
moves the state a long way without justification pays in the first term; one
that leaves a residual it could have removed pays in the second.

\begin{proposition}[Exact minimization]
$J$ decomposes over components, so
$\arg\min_{a} J(a)$ is obtained by minimizing each $J_{i}$ independently over
$\mathcal{C}(i) \cup \{\varnothing\}$.
\end{proposition}

\begin{remark}
This is worth stating plainly because it removes a whole apparatus. There is no
search here: the optimal assignment is a per-component $\arg\min$ over
$\kappa + 1$ precomputed numbers, and it is obtained in one pass over the cost
table. The decomposition is a property of the cost \eqref{eq:global}, not of
the implementation.
\end{remark}

\section{From assignment to analysis}

The assignment determines a radius per component, and therefore the predecessor
sets of the modified-Cholesky estimator \cite{bickel2008,nino2018}. For each
component $i$ with predecessors $P(i)$ given by its radius,
%
\begin{equation}
  \boldsymbol{\beta}_{i}
  = \big(\mathbf{X}_{P}^{\top}\mathbf{X}_{P} + a_{i}\mathbf{I}\big)^{-1}
    \mathbf{X}_{P}^{\top}\,\Delta\mathbf{X}_{i},
  \qquad a_{i} = \alpha\,\frac{\operatorname{tr}\big(\mathbf{X}_{P}^{\top}\mathbf{X}_{P}\big)}{|P(i)|},
  \qquad
  \Lmat_{ii} = 1, \quad \Lmat_{i,P(i)} = -\boldsymbol{\beta}_{i}^{\top},
\end{equation}
%
with $\Dmat_{ii}$ the inverse residual variance, and
$\widehat{\Bmat}^{-1} = \Lmat^{\top}\Dmat\Lmat$. The ridge is a fraction
$\alpha$ of the mean eigenvalue of each row's own Gram matrix rather than an
absolute constant. This matters here more than in the fixed-radius estimator:
the assignment produces rows with anywhere from zero to over a hundred
predecessors in the same cycle, and an absolute penalty would regularize the
short rows hard and the wide rows not at all. The relative form gives
$\alpha$ the same meaning in every row, every cycle and every configuration
compared below; $\alpha = 0.3$ throughout. The analysis solves
%
\begin{equation}
  \big(\widehat{\Bmat}^{-1} + \Hmat^{\top}\Rmat^{-1}\Hmat\big)\,
  \delta\mathbf{x} = \Hmat^{\top}\Rmat^{-1}\big(\mathbf{y} - \Hmat\xb\big).
\end{equation}
%
The precision is positive definite for any assignment, because it retains the
form $\Lmat^{\top}\Dmat\Lmat$; truncating an already-formed precision does not
preserve this, and in a companion measurement produced an indefinite matrix in
44 of 44 cases tested.

\begin{algorithm}[t]
\caption{One assimilation cycle}
\begin{algorithmic}[1]
\State propagate the ensemble; form anomalies $\Delta\mathbf{X}$
\State \textbf{for each component} grow its neighbourhood until it holds
       $\kappa$ observations
\State precompute $\sigma_{ii}$, $\sigma_{jj}$ and $\sigma_{ij}$ for every
       candidate pair
\State \textbf{for each component} take
       $a(i) = \arg\min J_{i}$ over $\mathcal{C}(i)\cup\{\varnothing\}$
\State read off the radius of every component from $a$
\State assemble $\widehat{\Bmat}^{-1}$ once and solve for the analysis
\end{algorithmic}
\end{algorithm}

\section{Numerical experiments}

The testbed is the 1.5-layer quasi-geostrophic model of Sakov and
Oke~\cite{sakov2008},
%
\begin{equation}
  q_t = -\psi_x - \varepsilon J(\psi,q) - A\triangle^{3}\psi + 2\pi\sin(2\pi y),
  \qquad q = \triangle\psi - F\psi,
\end{equation}
%
on the unit square with $F = 1600$, $\varepsilon = 10^{-5}$ and Dirichlet
boundaries, at $49 \times 49$ per field. The integration is carried on $q$ and
$\psi$ is recovered from it at every step, so $\psi$ is diagnostic and only $q$
is estimated. $q$ vanishes on the Dirichlet boundary for every state, so the
boundary components are removed from the assimilation altogether, as in
\cite{sakov2008}: they are not observed, not candidates, not predecessors,
and not in the error. The estimated state is the $47 \times 47$ interior,
$n = 2209$. Ensemble members and the truth are drawn from
snapshots of a long run, following \cite{sakov2008}. The dissipation is ten
times the model default, which is what gives a stationary regime: at the
default the norm of $q$ grows by a factor of fourteen between $t=500$ and
$t=8000$ and there is no climatology to sample.

\subsection{What the assignment does on a forecast}

Everything in this subsection is computed from ensemble anomalies and
observations before any analysis is performed, so it does not depend on
which filter follows. Table~\ref{tab:single_cycle} gives, for both
observation networks and for a climatological ensemble as well as one the
filter has produced, the cost of each rule and the structure of the
precision the greedy field induces against a uniform radius of the same
mean. Three things are stable across seeds, networks and cycles. The cost
separates the rules by orders of magnitude, and the optimal assignment
improves on the nearest observation by an order of magnitude. Roughly three
quarters of the components with a choice do not take their nearest
observation: the ensemble correlation, not the distance, decides. And a
third of the observations update no component at all; the ensemble variance
at those sites is one to two orders of magnitude below the variance at the
sites that are used, so the assignment is discarding observations where the
Kalman gain would be near zero anyway (Fig.~\ref{fig:discarded}).

\begin{figure}[t]\centering
\IfFileExists{../results/figures_paper/radius_fields_lattice_enkf-mc.pdf}{%
\includegraphics[width=\linewidth]{../results/figures_paper/radius_fields_lattice_enkf-mc.pdf}}{%
\fbox{\parbox{0.9\linewidth}{\centering\small\texttt{results/figures\_paper/radius\_fields\_lattice\_enkf-mc.pdf} --- run \texttt{make paper}}}}
\caption{The radius field of each rule inside the cycled EnKF-MC on the
lattice network, on one colormap, and the uniform field of the same mean as
the greedy one. Crosses mark observations the greedy assignment does not
use.}\label{fig:radius}
\end{figure}

\begin{figure}[t]\centering
\IfFileExists{../results/figures_paper/tessellation_lattice_enkf-mc.pdf}{%
\includegraphics[width=\linewidth]{../results/figures_paper/tessellation_lattice_enkf-mc.pdf}}{%
\fbox{\parbox{0.9\linewidth}{\centering\small\texttt{tessellation\_lattice\_enkf-mc.pdf}}}}
\caption{Which observation updates each component: the nearest-observation
tessellation against the optimal assignment, and the components where the
two differ.}\label{fig:tess}
\end{figure}

\begin{figure}[t]\centering
\IfFileExists{../results/figures_paper/discarded_obs_lattice_enkf-mc.png}{%
\includegraphics[width=\linewidth]{../results/figures_paper/discarded_obs_lattice_enkf-mc.png}}{%
\fbox{\parbox{0.9\linewidth}{\centering\small\texttt{discarded\_obs\_lattice\_enkf-mc.png}}}}
\caption{Ensemble spread of $q$ along the run with the observations that
update nothing marked, and the distribution of ensemble variance at used and
unused observation sites over all cycles.}\label{fig:discarded}
\end{figure}

\subsection{The cycled benchmark}

Four rules for the radius field are compared inside two filters: the
modified-Cholesky EnKF, where the field defines the predecessor sets of
$\widehat{\Bmat}^{-1}$, and the LETKF, where the same field defines the local
domain of each component. The uniform radius $r \in \{1, 2, 3, 4, 6, 8\}$ is
the standard scheme and the baseline; greedy is the method; nearest and
random are the two controls, one showing what distance alone achieves and
the other that a spread of radii is not by itself the point. Both filters
run on two networks of equal count: a lattice of spacing two, which shifts
by one grid point per cycle, and a random network redrawn every cycle.
$N = 40$, 60 cycles with the first 15 discarded, observations every 20 time
units with 5\% noise, inflation 1.15, $\kappa = 4$, three seeds. In the
EnKF-MC a component that abstains has a diagonal precision row; in the LETKF
it keeps its background. Divergence --- a forecast the model cannot integrate
--- is recorded per run and counted, not averaged.

Table~\ref{tab:rmse_main} reports the post-burn-in analysis error of $q$
over the interior; Table~\ref{tab:assignment} what the assignment did inside
each filter; Table~\ref{tab:timing} the cost per cycle;
Table~\ref{tab:sensitivity} the same comparison at $N = 20$, at
$\alpha = 0.1$, and on a lattice of spacing three, so that neither the
ensemble size, the regularization nor the density is left as a question. Figure~\ref{fig:rvsr} places the rules against the
uniform-radius sweep.

\begin{figure}[t]\centering
\IfFileExists{../results/figures_paper/rmse_vs_radius_main.pdf}{%
\includegraphics[width=\linewidth]{../results/figures_paper/rmse_vs_radius_main.pdf}}{%
\fbox{\parbox{0.9\linewidth}{\centering\small\texttt{rmse\_vs\_radius\_main.pdf}}}}
\caption{Post-burn-in analysis RMSE against the uniform radius (points,
median and range over seeds), with each assignment rule as a horizontal band
placed at its mean radius (diamond), for both filters and both networks.}
\label{fig:rvsr}
\end{figure}

\IfFileExists{tables_paper.tex}{\input{tables_paper.tex}}{%
\begin{center}\fbox{\parbox{0.9\linewidth}{\small\texttt{paper/tables\_paper.tex} is written by
\texttt{figures/make\_tables.py} after \texttt{make paper}; until then the
tables referenced above are absent.}}\end{center}}

\section{Discussion}

The contribution is a reparameterization. The radius stops being a free number
and becomes a consequence of a choice among observations that are actually
present, which is what makes it estimable: a choice among $\kappa+1$
alternatives is constrained in a way that a positive real number is not.

Because the cost decomposes, the optimum is exact and costs one pass over a
table of $n(\kappa+1)$ numbers. The whole expense of the method is in
precomputing that table from ensemble anomalies once per cycle, and in
assembling the precision from the assignment that results. Nothing is tuned:
$\kappa$ sets how far a neighbourhood grows before it stops, and every other
quantity --- the radius of each component, which observations are used and
which are not --- is read off the assignment.

\begin{thebibliography}{9}
\bibitem{bickel2008} P.~J. Bickel and E.~Levina.
  Regularized estimation of large covariance matrices.
  \emph{The Annals of Statistics}, 36(1):199--227, 2008.
\bibitem{nino2018} E.~D. Nino-Ruiz, A.~Sandu, and X.~Deng.
  An ensemble Kalman filter implementation based on modified Cholesky
  decomposition for inverse covariance matrix estimation.
  \emph{SIAM Journal on Scientific Computing}, 40(2):A867--A886, 2018.
\bibitem{sakov2008} P.~Sakov and P.~R. Oke.
  A deterministic formulation of the ensemble Kalman filter: an alternative to
  ensemble square root filters. \emph{Tellus A}, 60(2):361--371, 2008.
\end{thebibliography}

\end{document}
